\documentclass[a4paper,12pt]{article}

\usepackage{alltt, fancyvrb, url}
\usepackage{graphicx}
\usepackage[utf8]{inputenc}
\usepackage{float}
\usepackage{hyperref}
\usepackage[italian]{babel}
\usepackage[table]{xcolor}
\usepackage{array}
\usepackage{multirow}
\usepackage{titlesec}
\usepackage{listings}
\usepackage[italian]{cleveref}
\usepackage{acronym}
\usepackage{subcaption}

\renewcommand{\thesection}{\arabic{section}}

\sloppy
\setlength{\parindent}{0pt}

\title{\textbf{Regole Operative del Team}}
\author{}
\date{}

\begin{document}
\maketitle
\tableofcontents

\newpage

%\section{Problem solving} % meeting?

\section{Decision making}
Il processo di decision making seguirà l'approccio \textit{consultativo}, in cui il responsabile definito nella matrice
delle responsabilità prende la decisione, ma dopo aver raccolto informazioni e idee dai membri del team.

%\section{Conflict resolution} % ???
%\section{Consensus building} % ???
%\section{Brainstorming} % ???

\section{Team meetings}
Non verranno effettuate riunioni giornaliere, in quanto lo stato di avanzamento dei task è visibile a tutti i membri
del team tramite la Kanban board. Nel caso sorgano problemi durante lo sviluppo, il team pianificherà autonomamente
riunioni per discutere e risolvere i problemi.\\
In concomitanza con il raggiungimento delle milestone, si effettueranno project review meetings in cui si presenterà
lo stato del progetto ed il suo avanzamento, e se ne effettuerà una revisione. A queste riunioni parteciperanno i
project manager, i senior management, il committente, lo sponsor, e i membri principali del team di sviluppo.

\section{Gestione cambiamenti di scope}
In caso emerga la necessità di modificare lo scope del progetto, è necessario sottomettere una richiesta di cambiamento
dello scope al project manager, il quale esegue uno studio sull'impatto del cambiamento di scope sul progetto, ponendo
particolare attenzione al mantenimento dell'integrità concettuale, e durante l'intero processo potrà rifiutare la
richiesta. In caso il cambiamento possa essere applicato, la richiesta verrà approvata e verranno messe in campo le
misure necessarie per implementare il cambiamento.\\
Il form per sottomettere una richiesta di cambiamento dello scope deve contenere le seguenti informazioni:
\begin{itemize}
    \item nome del progetto;
    \item nominativo del richiedente;
    \item data della richiesta;
    \item descrizione del cambiamento;
    \item motivazione;
    \item spazio per l'eventuale approvazione (firma e data).
\end{itemize}

\section{Gestione delle comunicazioni}
Tutte le comunicazioni ufficiali devono essere riportate in forma scritta e devono essere condivise via mail con tutti
i soggetti interessati, i quali devono confermare la ricezione della comunicazione.\\
Al termine di ogni riunione, il verbale che è stato redatto deve essere condiviso con tutti i partecipanti, i quali
devono approvarlo.

%\section{Gestione delle risorse} % ???

\end{document}